\chapter{HCR-Inspired Pair Representation: Detailed Definition}
\label{app:hcr-details}

The statistical branch in Task~A represents each ordered variable pair $(U,V)$ by a fixed $40$-dimensional vector.
Its core construction is inspired by Hierarchical Correlation Reconstruction (HCR), where dependence is represented through coefficients of products of orthonormal basis functions~\cite{duda2018hcr,duda2018hcr-ts,duda2020hcr-income}.
The implementation extends this coefficient representation with fixed-width summaries describing dependence magnitude, marginal distributions, joint activation, data support, estimability and variable type.
This appendix gives the complete formula-level definition used by the final model.

Two layers must be kept distinct.
\begin{enumerate}[label=(\arabic*)]
\item \textbf{HCR inspiration.}
Dependence is encoded through coefficients of products of type-specific orthonormal basis functions, following the HCR principle of representing statistical dependence in an orthonormal expansion~\cite{duda2018hcr,duda2018hcr-ts,duda2020hcr-income}.
\item \textbf{Thesis-specific extensions.}
Fixed-width padding, derived dependence summaries, marginal descriptors, joint-activation descriptors, support and estimability features, and type metadata provide a common interface for heterogeneous binary, count and continuous variables.
\end{enumerate}
The resulting object is therefore called an \emph{HCR-inspired pair representation}.
It is a statistical dependence descriptor and does not identify causal direction.

\section{Overall $40$-dimensional representation}
\label{app:hcr-overall}

Define
\begin{equation}
h(U,V)\in\mathbb{R}^{40}.
\end{equation}
Let
\begin{equation}
I_{UV}=\{i:U_i\text{ and }V_i\text{ are both observed}\}
\end{equation}
and
\begin{equation}
n=|I_{UV}|,
\end{equation}
with $n_{\mathrm{train}}$ denoting the total number of training patients.
All bases and statistical descriptors are fitted from training patients only.

The slot composition is
\begin{center}
\begin{tabular}{@{}cl@{}}
\toprule
Slots & Content \\
\midrule
$0$--$15$ & padded HCR-inspired coefficient matrix \\
$16$--$19$ & derived dependence summaries \\
$20$--$23$ & marginal descriptors of $U$ \\
$24$--$27$ & marginal descriptors of $V$ \\
$28$--$31$ & joint activation / enrichment / rarity \\
$32$ & complete-case support \\
$33$ & estimability indicator \\
$34$--$36$ & one-hot type of $U$ \\
$37$--$39$ & one-hot type of $V$ \\
\bottomrule
\end{tabular}
\end{center}
Dimension check:
\begin{equation}
16+4+4+4+4+1+1+3+3=40.
\end{equation}

\begin{table}[htbp]
\centering
\small
\caption{Block layout of the $40$-dimensional HCR-inspired pair representation.}
\label{tab:app-hcr40-blocks}
\begin{tabular}{@{}ccl@{}}
\toprule
Slots & Dim.\ & Block \\
\midrule
$0$--$15$ & $16$ & padded coefficient matrix $A_{\mathrm{pad}}$ \\
$16$--$19$ & $4$ & derived dependence summaries \\
$20$--$23$ & $4$ & marginal descriptors $m(U)$ \\
$24$--$27$ & $4$ & marginal descriptors $m(V)$ \\
$28$--$31$ & $4$ & joint activation / enrichment / rarity \\
$32$ & $1$ & complete-case support $S$ \\
$33$ & $1$ & estimability indicator $I_{\mathrm{sup}}$ \\
$34$--$36$ & $3$ & one-hot type of $U$ \\
$37$--$39$ & $3$ & one-hot type of $V$ \\
\bottomrule
\end{tabular}
\end{table}

\section{HCR-inspired coefficient matrix}
\label{app:hcr-matrix}

For each variable construct type-specific non-constant orthonormal basis functions
\begin{equation}
\phi^{(U)}_{j},\quad j=1,\ldots,d_{U},
\qquad
\phi^{(V)}_{k},\quad k=1,\ldots,d_{V}.
\end{equation}
Define
\begin{equation}
\label{eq:app-ajk}
a_{jk}
=
\frac{1}{n}
\sum_{i\in I_{UV}}
\phi^{(U)}_{j}(U_i)\,
\phi^{(V)}_{k}(V_i).
\end{equation}
Equivalently, with design matrices $\Phi_U$ and $\Phi_V$ whose rows are the evaluated basis functions on complete cases,
\begin{equation}
A_{UV}=\frac{1}{n}\Phi_U^{\mathsf{T}}\Phi_V.
\end{equation}
Zero-pad the active block to
\begin{equation}
A_{\mathrm{pad}}\in\mathbb{R}^{4\times 4}
\end{equation}
and store it row-wise (row-major) in slots~$0$--$15$:
\begin{equation}
A_{\mathrm{pad}}
=
\begin{bmatrix}
a_{11} & a_{12} & a_{13} & a_{14}\\
a_{21} & a_{22} & a_{23} & a_{24}\\
a_{31} & a_{32} & a_{33} & a_{34}\\
a_{41} & a_{42} & a_{43} & a_{44}
\end{bmatrix}.
\end{equation}
Active block sizes by pair type are
\begin{center}
\begin{tabular}{@{}ll@{}}
\toprule
Pair type & Active block \\
\midrule
binary $\times$ binary & $1\times 1$ \\
binary $\times$ continuous & $1\times 4$ \\
binary $\times$ count & $1\times 4$ \\
count $\times$ count & up to $4\times 4$ \\
count $\times$ continuous & up to $4\times 4$ \\
continuous $\times$ continuous & up to $4\times 4$ \\
\bottomrule
\end{tabular}
\end{center}
The coefficient $a_{11}$ describes the lowest-order dependence mode; higher-order coefficients allow the representation to capture additional nonlinear structure rather than reducing every pair to one correlation number.
Individual higher-order coefficients are not assigned unique causal or clinical meanings.

\section{Binary basis and connection to Pearson $\phi$}
\label{app:hcr-binary}

For $B\in\{0,1\}$, with training prevalence $p_{B}=P(B=1)$, define the standardized non-constant basis function
\begin{equation}
\label{eq:app-phi-binary}
\phi^{(B)}_{1}(B)
=
\frac{B-p_{B}}{\sqrt{p_{B}(1-p_{B})}}.
\end{equation}
For a binary--binary pair,
\begin{equation}
\label{eq:app-a11-phi}
a_{11}
=
\frac{p_{11}-p_{U}p_{V}}{\sqrt{p_{U}(1-p_{U})\,p_{V}(1-p_{V})}},
\end{equation}
which is equivalent to the Pearson $\phi$ coefficient.
For binary pairs, $a_{11}$ therefore reduces to Pearson $\phi$.
This is the conceptual link between the simple binary association used elsewhere and the more general HCR-inspired coefficient representation.

\section{Continuous variables}
\label{app:hcr-continuous}

Continuous variables are transformed using the empirical training distribution.
The training rank / pseudo-observation construction is
\begin{equation}
\label{eq:app-ecdf}
u_{i}
=
\mathrm{clip}\!\left(
\frac{\tfrac{1}{2}(\ell_{i}+r_{i})+0.5}{N_{\mathrm{train}}+1},
10^{-6},\,1-10^{-6}
\right),
\end{equation}
where $\ell_{i}$ and $r_{i}$ are the left and right rank positions of tied observations in the sorted training sample.
The transformed value is expanded in shifted orthonormal Legendre functions
\begin{equation}
\label{eq:app-legendre}
\psi_{j}(u)=\sqrt{2j+1}\,P_{j}(2u-1),\qquad j=1,\ldots,4,
\end{equation}
where $P_{j}$ is the Legendre polynomial of order $j$.
The rank transformation removes dependence on the original marginal scale, while the orthonormal polynomial expansion provides several low-order modes with which joint dependence can be represented.
Legendre polynomials are the chosen orthonormal basis in this implementation; they are not themselves the HCR method.

\section{Count variables}
\label{app:hcr-count}

Counts are capped using the training $99.5$th percentile,
\begin{equation}
\label{eq:app-count-cap}
C^{c}=\min(C,c_{\max}),\qquad c_{\max}=q_{0.995}.
\end{equation}
Repeated integer values are handled with deterministic distributional jitter,
\begin{equation}
\label{eq:app-jitter}
u_{i}=F_{C}(c_{i}^{-})+\xi_{i}\,P(C=c_{i}),\qquad \xi_{i}\in(0,1),
\end{equation}
where $\xi_{i}$ is deterministic and reproducible.
The transformed values are then passed through the same four shifted Legendre basis functions used for continuous variables.
Count variables are \emph{not} simply binarized in the final model.

\section{Four derived dependence summaries (slots~$16$--$19$)}
\label{app:hcr-summaries}

The following four deterministic summaries of the active coefficient matrix are stored in slots~$16$--$19$.

For slot~$16$ (total energy),
\begin{equation}
D_{\mathrm{tot}}=\|A\|_{F}^{2}=\sum_{j,k}a_{jk}^{2}.
\end{equation}
For slot~$17$ (mean energy),
\begin{equation}
D_{\mathrm{mean}}=\frac{\|A\|_{F}^{2}}{d_{U}\,d_{V}}.
\end{equation}
For slot~$18$ (low-order energy fraction),
\begin{equation}
R_{\mathrm{low}}
=
\frac{\|A_{1:2,1:2}\|_{F}^{2}}{\|A\|_{F}^{2}+\varepsilon},
\end{equation}
with $\varepsilon=10^{-12}$ in the implementation.
For slot~$19$ (strongest coefficient),
\begin{equation}
M_{\max}=\max_{j,k}\lvert a_{jk}\rvert.
\end{equation}

Interpretation:
$D_{\mathrm{tot}}$ is overall dependence magnitude;
$D_{\mathrm{mean}}$ normalizes that magnitude by the number of active coefficients;
$R_{\mathrm{low}}$ is the fraction of dependence concentrated in low-order modes;
$M_{\max}$ is the strength of the strongest individual dependence mode.
These four values are deterministic summaries of $A$, not independent statistical estimators.

A paired ablation that removed slots~$16$--$19$ (reducing the vector from $40$ to $36$ dimensions) lowered validation Average Precision (AP) by approximately $0.008$ on average, which supports retaining the summaries while treating them as derived rather than independent evidence.

\section{Type-specific marginal descriptors (slots~$20$--$27$)}
\label{app:hcr-marginals}

Each variable contributes a four-dimensional marginal vector
\begin{equation}
m(U)\in\mathbb{R}^{4},\qquad m(V)\in\mathbb{R}^{4},
\end{equation}
with slots~$20$--$23$ storing $m(U)$ and slots~$24$--$27$ storing $m(V)$.

\subsection*{Binary}
\begin{enumerate}[label=\arabic*.]
\item prevalence $p_{1}=P(B=1)$;
\item normalized binary entropy
\begin{equation}
H_{2}(p_{1})
=
\frac{-p_{1}\ln p_{1}-(1-p_{1})\ln(1-p_{1})}{\ln 2};
\end{equation}
\item bounded imbalance $\tanh\bigl(\mathrm{logit}(p_{1})/4\bigr)$, where
\begin{equation}
\mathrm{logit}(p)=\ln\frac{p}{1-p};
\end{equation}
\item positive-state rarity $(n_{1}+1)^{-1/2}$.
\end{enumerate}

\subsection*{Count}
Using the capped count $C^{c}$:
\begin{enumerate}[label=\arabic*.]
\item zero mass $P(C^{c}=0)$;
\item scaled count intensity
\begin{equation}
\frac{\mathbb{E}[\log(1+C^{c})]}{\log(1+c_{\max})};
\end{equation}
\item normalized empirical count entropy $H(\hat{\pi})/\log K_{\mathrm{obs}}$, where
\begin{equation}
K_{\mathrm{obs}}=\max\bigl(|\mathrm{unique}(C^{c})|,2\bigr)
\end{equation}
and
\begin{equation}
H(\hat{\pi})=-\sum_{r}\hat{\pi}_{r}\log\hat{\pi}_{r};
\end{equation}
\item bounded dispersion descriptor
\begin{equation}
\tanh\!\left(
\log\frac{\mathrm{Var}(C^{c})+\varepsilon}{\mathbb{E}[C^{c}]+\varepsilon}
\right).
\end{equation}
\end{enumerate}

\subsection*{Continuous}
Define the robust standardized value
\begin{equation}
\label{eq:app-zr}
z_{r}
=
\frac{x-\mathrm{median}(x)}{1.4826\,\mathrm{MAD}(x)+\varepsilon},
\end{equation}
with
\begin{equation}
\mathrm{MAD}(x)=\mathrm{median}\bigl(\lvert x-\mathrm{median}(x)\rvert\bigr).
\end{equation}
Let
\begin{equation}
z_{0}
=
\frac{z_{r}-\mathbb{E}[z_{r}]}{\mathrm{std}(z_{r})+\varepsilon},
\end{equation}
using the population standard deviation over finite training observations, and set
\begin{equation}
\gamma_{1}=\mathbb{E}[z_{0}^{3}],\qquad \gamma_{2}=\mathbb{E}[z_{0}^{4}].
\end{equation}
The four continuous descriptors are
\begin{enumerate}[label=\arabic*.]
\item bounded skewness $\tanh(\gamma_{1}/3)$;
\item bounded excess kurtosis / tail-weight descriptor $\tanh((\gamma_{2}-3)/10)$;
\item extreme tail mass $P(\lvert z_{r}\rvert>1.5)$;
\item signed tail asymmetry $P(z_{r}>1.5)-P(z_{r}<-1.5)$.
\end{enumerate}
Skewness and kurtosis use $z_{0}$; the tail descriptors continue to use $z_{r}$.
The hyperbolic tangent is used only as a bounded monotonic transform, not as an additional statistical measure.

\section{Joint activation and rarity (slots~$28$--$31$)}
\label{app:hcr-joint}

Define a type-specific activation indicator
\begin{equation}
\label{eq:app-act}
\mathrm{act}(X)
=
\begin{cases}
\mathbf{1}[X\ge 0.5], & X\text{ binary},\\
\mathbf{1}[X>0], & X\text{ count},\\
\mathbf{1}[\lvert z_{r}\rvert>1.5], & X\text{ continuous}.
\end{cases}
\end{equation}
Marginal activation rates are
\begin{equation}
p_{+}(U)=\frac{1}{n}\sum_{i}\mathrm{act}(U_{i}),
\qquad
p_{+}(V)=\frac{1}{n}\sum_{i}\mathrm{act}(V_{i}),
\end{equation}
and joint activation is
\begin{equation}
p_{++}
=
\frac{1}{n}
\sum_{i}
\mathbf{1}\bigl[\mathrm{act}(U_{i})=1\wedge\mathrm{act}(V_{i})=1\bigr].
\end{equation}
\begin{itemize}
\item slot~$28$: joint activation rate $p_{++}$;
\item slot~$29$: excess joint activation over independence $\Delta=p_{++}-p_{+}(U)p_{+}(V)$;
\item slot~$30$: log-lift
\begin{equation}
L=\log\frac{p_{++}+\varepsilon}{p_{+}(U)p_{+}(V)+\varepsilon};
\end{equation}
\item slot~$31$: joint-rarity descriptor
\begin{equation}
n_{++}=\sum_{i}\mathbf{1}\bigl[\mathrm{act}(U_{i})=1\wedge\mathrm{act}(V_{i})=1\bigr],
\qquad
R_{++}=(n_{++}+1)^{-1/2}.
\end{equation}
\end{itemize}
The quantity $R_{++}$ is an engineered rarity descriptor, not a classical rarity estimator.

\section{Support and estimability (slots~$32$--$33$)}
\label{app:hcr-support}

For slot~$32$,
\begin{equation}
S=\frac{n_{\mathrm{complete}}}{n_{\mathrm{train}}},
\end{equation}
which communicates how much of the training population supports the pair estimate.
For slot~$33$, $I_{\mathrm{sup}}\in\{0,1\}$.
The estimability indicator distinguishes a genuinely small or near-zero dependence estimate from an unavailable estimate that was zero-filled because support was insufficient.

In the frozen implementation a pair is unsupported if
\begin{equation}
n_{\mathrm{complete}}<50,
\end{equation}
and, for binary--binary pairs, if the minimum of the four level counts
\begin{equation}
\lvert\{U=0\}\rvert,\;
\lvert\{U=1\}\rvert,\;
\lvert\{V=0\}\rvert,\;
\lvert\{V=1\}\rvert
\end{equation}
among complete cases is smaller than $5$.
When unsupported, the coefficient matrix is zeroed while enrichment and type metadata may still be filled.

\section{Variable-type metadata (slots~$34$--$39$)}
\label{app:hcr-types}

slots~$34$--$36$ store
\begin{equation}
\mathrm{type}(U)=\mathrm{one\text{-}hot}(\mathrm{binary},\mathrm{count},\mathrm{continuous}),
\end{equation}
and slots~$37$--$39$ store the analogous one-hot vector for $V$.
These indicators are needed because zero-padding has different meanings for different pair types:
for a binary--binary pair most of the $4\times 4$ coefficient matrix is structurally absent, whereas for a continuous--continuous pair a zero coefficient may represent a genuinely small estimated dependence mode.

\section{Final composition}
\label{app:hcr-compose}

The complete vector is
\begin{equation}
\begin{aligned}
h(U,V)
=
\bigl[
&\mathrm{vec}(A_{\mathrm{pad}})
\parallel
D_{\mathrm{tot}}
\parallel
D_{\mathrm{mean}}
\parallel
R_{\mathrm{low}}
\parallel
M_{\max}
\\
&\parallel
m(U)
\parallel
m(V)
\parallel
p_{++}
\parallel
\Delta
\parallel
L
\parallel
R_{++}
\\
&\parallel
S
\parallel
I_{\mathrm{sup}}
\parallel
t_{U}
\parallel
t_{V}
\bigr]
\in\mathbb{R}^{40},
\end{aligned}
\end{equation}
where $\parallel$ denotes concatenation and $\mathrm{vec}$ is row-major flattening.
Dimension check:
\begin{equation}
16+4+4+4+4+1+1+3+3=40,
\end{equation}
where the four after the two marginal blocks are the joint-activity features.

\section{Connection to the Task~A motif}
\label{app:hcr-motif}

For candidate $A\to G$ with local context node $Z$, the statistical branch constructs independently
\begin{equation}
h_{AZ}=h(A,Z),\qquad
h_{AG}=h(A,G),\qquad
h_{ZG}=h(Z,G),
\end{equation}
each in $\mathbb{R}^{40}$.
These three vectors are not immediately concatenated into the fusion input.
Each role is first passed through its own $40\to 16\to 8$ role encoder, and absent motif roles are masked,
\begin{equation}
e_{r}=m_{r}\,E_{r}(h_{r}),
\end{equation}
yielding
\begin{equation}
g_{\mathrm{stat}}
=
\bigl[e_{AZ}\parallel e_{AG}\parallel e_{ZG}\bigr]
\in\mathbb{R}^{24}.
\end{equation}

\section{Source and inspiration note}
\label{app:hcr-source}

The coefficient construction is inspired by Hierarchical Correlation Reconstruction: dependence is encoded through expectations of products of orthonormal basis functions~\cite{duda2018hcr,duda2018hcr-ts,duda2020hcr-income}.
The fixed $40$-dimensional representation used here is an application-specific extension of that principle rather than a direct reproduction of a canonical HCR feature vector.
In particular, the marginal, joint-activation, support, rarity and type descriptors were introduced to provide a common interface for heterogeneous binary, count and continuous variables and are not attributed to the original HCR formulation.
